% ===================================================================
% Section 4: Theoretical properties
% ===================================================================
\section{Theoretical properties}
\label{sec:theory}

This section develops the structural theory of program \eqref{eq:objective}. We
first fix standing assumptions, then establish convexity, existence, uniqueness,
first-order conditions, exact reductions to the classical portfolios, monotone
comparative statics in the penalty strengths, limiting solutions, and a stability
bound.

\subsection{Standing assumptions}
\label{sec:assumptions}

\begin{assumption}[Risk model]
\label{as:sigma}
$\Sighat\in\Sym^{N}$ is symmetric positive semidefinite. Where strict convexity
is invoked we assume in addition that either $\Sighat\succ0$ or $\lambda_2>0$.
\end{assumption}

\begin{assumption}[Design matrix]
\label{as:V}
$V=[\vv_1,\dots,\vv_N]\in\R^{K\times N}$ with $K\le N$ and $\rank V=K$.
\end{assumption}

\begin{assumption}[Normalization]
\label{as:norm}
$V$ is scaled so that $\tfrac1N\tr\!\big(VV^\top\big)
=\tfrac1N\sum_{i=1}^{N}\lVert\vv_i\rVert_2^{2}=1$.
\end{assumption}

\begin{assumption}[Feasible set]
\label{as:feasible}
$\Wset$ in \eqref{eq:feasible} is nonempty and, where compactness is invoked, the
box bounds are finite (in particular $\Wset$ is compact in the long-only case
$\ell=0$, being a subset of the unit simplex). Slater's condition holds:
$\relint\Wset\neq\emptyset$, i.e.\ there exists a feasible $\w$ with
$\ell_i<w_i<u_i$ for all $i$ and $\sum_{i\in g}w_i<c_g$ for all $g$.
\end{assumption}

\begin{assumption}[Effective domain]
\label{as:domain}
When $\gA>0$ or $\gD>0$ we restrict attention to
$\dom F=\{\w:\M\succ0\}$, on which $F$ is real valued, and set $F=+\infty$
elsewhere; by Corollary~\ref{cor:spanning} this requires
$\{\vv_i:w_i>0\}$ to span $\R^{K}$, and we assume
$\Wset\cap\dom F\neq\emptyset$.
\end{assumption}

\begin{remark}
Assumption~\ref{as:domain} is not vacuous. In the square case $K=N$ it reduces to
$\w>0$, so the effective domain of the A- and D-penalized problem is the open
simplex; in the factor case $K<N$ it is the weaker requirement that the assets
held span factor space, which permits some $w_i=0$ provided at least $K$ assets
with linearly independent loadings are held.
\end{remark}

\subsection{Derivatives}
\label{sec:gradients}

Since $\partial\M/\partial w_i=\vv_i\vv_i^\top$, the chain rule gives closed forms
for every term. Define the classical \emph{variance function} of the design,
\begin{equation}
  d_i(\w)\;=\;\vv_i^\top \M^{-1}\vv_i,\qquad i=1,\dots,N ,
  \label{eq:varfun}
\end{equation}
which in the square case $K=N$ reduces to $d_i(\w)=1/w_i$.

\begin{lemma}[Gradients]
\label{lem:gradients}
On $\dom F$, with $\uu^\star$ a unit eigenvector associated with $\lammin(\M)$,
\begin{align}
  \nabla_{\w}\big(\w^\top\Sighat\w\big) &= 2\Sighat\w, &
  \nabla_{\w}\big(\lambda_2\lVert\w\rVert_2^2\big) &= 2\lambda_2\w,
  \label{eq:grad1}\\
  \frac{\partial}{\partial w_i}\,\phiA(\M) &= -\,\vv_i^\top \M^{-2}\vv_i
  \;=\;-\big\lVert \M^{-1}\vv_i\big\rVert_2^{2}, &
  \frac{\partial}{\partial w_i}\,\phiD(\M) &= -\,d_i(\w),
  \label{eq:grad2}\\
  \frac{\partial}{\partial w_i}\,\phiE(\M) &= -\big(\vv_i^\top\uu^\star\big)^{2}.
  & &
  \label{eq:grad3}
\end{align}
The expression \eqref{eq:grad3} is a gradient when $\lammin(\M)$ is a simple
eigenvalue and a subgradient otherwise.
\end{lemma}

\begin{proof}
\eqref{eq:grad1} is standard. For $\phiA$, differentiating
$\M\M^{-1}=I$ gives $\partial\M^{-1}/\partial w_i=-\M^{-1}(\partial\M/\partial w_i)\M^{-1}
=-\M^{-1}\vv_i\vv_i^\top\M^{-1}$, so
$\partial\tr(\M^{-1})/\partial w_i=-\tr(\M^{-1}\vv_i\vv_i^\top\M^{-1})
=-\vv_i^\top\M^{-2}\vv_i$, which equals $-\lVert\M^{-1}\vv_i\rVert_2^2$ by
symmetry of $\M$. For $\phiD$, Jacobi's formula gives
$\partial\log\det\M/\partial w_i=\tr(\M^{-1}\vv_i\vv_i^\top)=\vv_i^\top\M^{-1}\vv_i
=d_i(\w)$, and $\phiD=-\log\det\M$. For $\phiE$, first-order eigenvalue
perturbation gives
$\partial\lammin/\partial w_i=\uu^{\star\top}(\partial\M/\partial w_i)\uu^\star
=(\vv_i^\top\uu^\star)^2$ when the eigenvalue is simple; when it is repeated,
$\lammin$ is the pointwise minimum \eqref{eq:lammin-var} of linear functions and
$(\vv_i^\top\uu^\star)^2$ is an element of its superdifferential, so
$-(\vv_i^\top\uu^\star)^2$ is a subgradient of $\phiE$.
\end{proof}

\begin{remark}[Connection to the equivalence theory]
The appearance of the variance function \eqref{eq:varfun} as the negative
gradient of the D-criterion is the entry point to the equivalence theory of
optimal design: at an unpenalized D-optimal design the Kiefer--Wolfowitz theorem
\cite{kiefer1960} characterizes optimality by $\max_i d_i(\w)=K$. A
corresponding characterization for the penalized program \eqref{eq:objective},
via vertex directional derivatives on the simplex, is a natural development and
is treated separately.
\end{remark}

\subsection{Convexity}

\begin{proposition}[Convexity]
\label{prop:convex}
Under Assumptions~\ref{as:sigma}--\ref{as:domain}, and for
$\lambda_2,\gA,\gD,\gE\ge0$, the objective $F$ in \eqref{eq:objective} is a
proper closed convex function on $\dom F$, and $\Wset$ is a convex polyhedron.
Program \eqref{eq:objective} is therefore a convex program.
\end{proposition}

\begin{proof}
We treat the terms separately; a nonnegative combination of convex functions is
convex.

\emph{(i) Risk.} $\w\mapsto\w^\top\Sighat\w$ has Hessian $2\Sighat\succeq0$ by
Assumption~\ref{as:sigma}, hence is convex.

\emph{(ii) Ridge.} $\lambda_2\lVert\w\rVert_2^2$ has Hessian
$2\lambda_2I\succeq0$ for $\lambda_2\ge0$.

\emph{(iii) A-term.} By Lemma~\ref{lem:trinv} in the appendix, the map
$X\mapsto\tr(X^{-1})$ is convex on $\Sym^{K}_{++}$. The map $\w\mapsto\M$ is
linear, and convexity is preserved under composition with an affine map, so
$\w\mapsto\phiA(\M)$ is convex on $\dom F$.

\emph{(iv) D-term.} By Lemma~\ref{lem:logdet} in the appendix, $X\mapsto\log\det X$
is concave on $\Sym^{K}_{++}$, so $X\mapsto-\log\det X$ is convex; composing with
the linear map $\w\mapsto\M$ gives convexity of $\w\mapsto\phiD(\M)$.

\emph{(v) E-term.} By \eqref{eq:lammin-var}, $\w\mapsto\lammin(\M)$ is a pointwise
infimum over $\uu$ of the functions $\w\mapsto\sum_iw_i(\vv_i^\top\uu)^2$, each of
which is linear in $\w$. A pointwise infimum of linear functions is concave,
hence $\phiE(\M)=-\lammin(\M)$ is convex in $\w$.

\emph{Closedness and properness.} Each term is continuous on the open set
$\{\w:\M\succ0\}$, and $\phiA,\phiD\to+\infty$ as $\M$ approaches the boundary of
the positive definite cone, so extending $F$ by $+\infty$ off $\dom F$ yields a
lower semicontinuous, i.e.\ closed, function; it is proper by
Assumption~\ref{as:domain}.

\emph{Feasible set.} $\Wset$ is the intersection of the hyperplane
$\{\one^\top\w=1\}$, the box $\{\ell\le\w\le u\}$, and finitely many half-spaces
$\{\sum_{i\in g}w_i\le c_g\}$, each convex; hence $\Wset$ is a convex polyhedron.
\end{proof}

\subsection{Existence, uniqueness, and interiority}

\begin{theorem}[Existence]
\label{thm:existence}
Suppose Assumptions~\ref{as:sigma}--\ref{as:domain} hold and that either
\emph{(a)} $\Wset$ is compact, or \emph{(b)} $\lambda_2>0$. Then
\eqref{eq:objective} attains its minimum on $\Wset$.
\end{theorem}

\begin{proof}
By Proposition~\ref{prop:convex}, $F$ is closed, proper and convex, and
$\Wset\cap\dom F\neq\emptyset$.

\emph{(a)} If $\Wset$ is compact, let $(\w^{(m)})$ be a minimizing sequence in
$\Wset\cap\dom F$. By compactness it has a subsequence converging to some
$\w^{\infty}\in\Wset$. Lower semicontinuity of $F$ gives
$F(\w^{\infty})\le\liminf_mF(\w^{(m)})=\inf_{\Wset}F$, and since
$\inf_{\Wset}F<+\infty$ we get $\w^{\infty}\in\dom F$ and $\w^{\infty}$ is a
minimizer.

\emph{(b)} If $\lambda_2>0$ we show $F$ is coercive on $\Wset$, so that its
sublevel sets are bounded and case~(a) applies to a compact sublevel set. Indeed
$\w^\top\Sighat\w\ge0$ and $\phiA(\M)\ge0$; moreover
$\lammin(\M)\le\lammax(\M)\le\sum_iw_i\lVert\vv_i\rVert_2^2
\le\big(\max_i\lVert\vv_i\rVert_2^2\big)\lVert\w\rVert_1$, so
$\gE\phiE(\M)\ge-c_1\lVert\w\rVert_1$ for a constant $c_1$; and by
\eqref{eq:cb}, $\det\M$ is a polynomial of degree $K$ in $\w$, so
$\gD\phiD(\M)\ge-c_2-c_3\log(1+\lVert\w\rVert_2)$ for constants $c_2,c_3$. Hence
\[
  F(\w)\;\ge\;\lambda_2\lVert\w\rVert_2^2-c_1\lVert\w\rVert_1
  -c_3\log\big(1+\lVert\w\rVert_2\big)-c_2\;\longrightarrow\;+\infty
\]
as $\lVert\w\rVert_2\to\infty$, since the quadratic term dominates both a linear
and a logarithmic term.
\end{proof}

\begin{theorem}[Uniqueness]
\label{thm:uniqueness}
If any one of the following holds, $F$ is strictly convex on $\Wset\cap\dom F$ and
the minimizer of \eqref{eq:objective} is unique:
\begin{enumerate}[label=(\roman*),leftmargin=2.0em,itemsep=0.15em]
  \item $\Sighat\succ0$;
  \item $\lambda_2>0$;
  \item $\gD>0$ (or $\gA>0$) and $K=N$;
  \item $\gD>0$ (or $\gA>0$), $K<N$, and the map $\w\mapsto\M$ is injective.
\end{enumerate}
\end{theorem}

\begin{proof}
In cases (i) and (ii) the quadratic part $\w^\top(\Sighat+\lambda_2I)\w$ has
positive definite Hessian, hence is strictly convex, and adding the convex
remaining terms preserves strict convexity. In case (iii), by
Corollary~\ref{cor:KN} the D-term is $-\gD\sum_i\log w_i$, whose Hessian
$\gD\diag(w_i^{-2})$ is positive definite on $\w>0$; similarly the A-term has
Hessian $2\gA\diag(G_{ii}w_i^{-3})\succ0$ when $G_{ii}>0$. In case (iv),
$X\mapsto-\log\det X$ is strictly convex on $\Sym^{K}_{++}$, and the composition
of a strictly convex function with an \emph{injective} linear map is strictly
convex. Strict convexity on a convex set implies at most one minimizer, which
together with Theorem~\ref{thm:existence} gives existence and uniqueness.
\end{proof}

\begin{remark}[When is $\w\mapsto\M$ injective?]
\label{rem:injective}
The map $\w\mapsto\sum_iw_i\vv_i\vv_i^\top$ is injective precisely when the
rank-one matrices $\{\vv_i\vv_i^\top\}_{i=1}^{N}$ are linearly independent in
$\Sym^{K}$, which has dimension $K(K+1)/2$. A necessary condition is therefore
$N\le K(K+1)/2$. In a typical factor application with $K$ small and $N$ large
this \emph{fails}, so in the rank-deficient case the design penalties alone do
not deliver uniqueness and one should retain $\lambda_2>0$ (or $\Sighat\succ0$).
This is a concrete reason to keep the ridge term in the formulation rather than
regard it as redundant.
\end{remark}

\begin{proposition}[Interiority]
\label{prop:interior}
Let $K=N$ and $\gD>0$ (or $\gA>0$), and let $\Wset$ be contained in the unit
simplex. Then every minimizer $\w^\star$ of \eqref{eq:objective} satisfies
$w_i^\star>0$ for all $i$; that is, the nonnegativity constraints are inactive
and their multipliers vanish.
\end{proposition}

\begin{proof}
By Corollary~\ref{cor:KN} the D- and A-terms equal $-\gD\sum_i\log w_i$ and
$\gA\sum_iG_{ii}/w_i$ respectively, each of which tends to $+\infty$ as any
$w_i\downarrow0$, while the remaining terms stay bounded on the compact simplex.
Hence $F(\w)\to+\infty$ on any sequence approaching a face $\{w_i=0\}$, so no such
point can be a minimizer, and by Theorem~\ref{thm:existence} a minimizer exists in
the relative interior.
\end{proof}

Proposition~\ref{prop:interior} has a practical reading: with a strictly positive
D- or A-penalty the long-only constraint never binds, so the barrier
\emph{implements} the no-short-selling restriction rather than merely coexisting
with it.

\subsection{First-order conditions}

\begin{theorem}[KKT conditions: necessity and sufficiency]
\label{thm:kkt}
Let Assumptions~\ref{as:sigma}--\ref{as:domain} hold, with Slater's condition as
in Assumption~\ref{as:feasible}. Then $\w^\star\in\Wset\cap\dom F$ is a global
minimizer of \eqref{eq:objective} if and only if there exist a scalar $\nu$,
vectors $\underline{\alpha},\overline{\alpha}\in\R^{N}_{+}$, and scalars
$\mu_g\ge0$ such that
\begin{equation}
  2\Sighat\w^\star+2\lambda_2\w^\star
  -\gA\,\mathbf d_{\mathrm{A}}
  -\gD\,\mathbf d(\w^\star)
  -\gE\big(V^\top\uu^\star\big)^{\odot2}
  \;=\;\nu\one+\overline{\alpha}-\underline{\alpha}
  +\sum_g\mu_g\one_g ,
  \label{eq:kkt}
\end{equation}
where $(\mathbf d_{\mathrm{A}})_i=\vv_i^\top(\M^\star)^{-2}\vv_i$,
$\mathbf d(\w^\star)=(d_1(\w^\star),\dots,d_N(\w^\star))^\top$ is the variance
function \eqref{eq:varfun}, $(\cdot)^{\odot2}$ is entrywise squaring, and
$\one_g$ is the indicator of group $g$; together with primal feasibility
$\w^\star\in\Wset$, dual feasibility
$\underline{\alpha},\overline{\alpha},\mu_g\ge0$, and complementary slackness
\[
  \underline{\alpha}_i\big(w_i^\star-\ell_i\big)=0,\qquad
  \overline{\alpha}_i\big(u_i-w_i^\star\big)=0,\qquad
  \mu_g\Big(\textstyle\sum_{i\in g}w_i^\star-c_g\Big)=0 .
\]
\end{theorem}

\begin{proof}
All constraints are affine and Slater's condition holds, so the KKT conditions
are necessary for optimality. Conversely, since $F$ is convex
(Proposition~\ref{prop:convex}) and the constraints are affine, the KKT
conditions are sufficient: if $\w^\star$ satisfies them, then for any feasible
$\w$ the gradient inequality gives
$F(\w)\ge F(\w^\star)+\nabla F(\w^\star)^\top(\w-\w^\star)$, and substituting
\eqref{eq:kkt} together with feasibility and complementary slackness shows
$\nabla F(\w^\star)^\top(\w-\w^\star)\ge0$, whence $F(\w)\ge F(\w^\star)$. The
stationarity equation \eqref{eq:kkt} is Lemma~\ref{lem:gradients} substituted into
the stationarity of the Lagrangian; when $\lammin(\M^\star)$ is repeated,
$\uu^\star$ is replaced by an appropriate element of the superdifferential.
\end{proof}

\begin{corollary}[Global optimality]
\label{cor:global}
Because \eqref{eq:objective} is a convex program, every local minimizer is a
global minimizer, and the set of minimizers is convex and closed. If in addition
any condition of Theorem~\ref{thm:uniqueness} holds, this set is a singleton.
\end{corollary}

\begin{proof}
Let $\w^{\circ}$ be a local minimizer and $\w\in\Wset\cap\dom F$ arbitrary. For
$t\in(0,1]$ small, convexity of $\Wset\cap\dom F$ gives
$\w_t=(1-t)\w^{\circ}+t\w$ feasible and, for $t$ small enough,
$F(\w^{\circ})\le F(\w_t)\le(1-t)F(\w^{\circ})+tF(\w)$, whence
$tF(\w^{\circ})\le tF(\w)$ and so $F(\w^{\circ})\le F(\w)$. Thus $\w^{\circ}$ is
global. The minimizer set is the intersection of the convex set
$\Wset\cap\dom F$ with the sublevel set $\{F\le\min F\}$, which is convex and
closed since $F$ is convex and closed.
\end{proof}

\begin{remark}[Convexity versus the solver]
\label{rem:solver}
Corollary~\ref{cor:global} is a statement about the \emph{problem}, and it is
what licenses calling a computed stationary point a global optimum. It should not
be conflated with a property of any particular algorithm. Our implementation uses
sequential quadratic programming (SLSQP), which is a general-purpose nonlinear
constrained solver and does \emph{not} by itself certify global optimality; what
certifies it here is convexity. In practice we additionally verify
\eqref{eq:kkt} at the returned point and check the exact reductions of
Lemma~\ref{lem:reductions}. The same program can be posed for a conic solver---the
risk, ridge and A-terms admit semidefinite representations and the E-term is an
eigenvalue constraint---which returns a duality-gap certificate; we note this as
an alternative in Section~\ref{sec:computation}.
\end{remark}

\subsection{Exact reductions}

\begin{lemma}[Reductions]
\label{lem:reductions}
Let $\Sighat\succ0$, impose only the budget constraint $\one^\top\w=1$, and allow
short sales. Then:
\begin{enumerate}[label=(\roman*),leftmargin=2.0em,itemsep=0.15em]
  \item if $\lambda_2=\gA=\gD=\gE=0$, the unique solution of
    \eqref{eq:objective} is the GMV portfolio
    $\w^\star=\Sighat^{-1}\one/(\one^\top\Sighat^{-1}\one)$;
  \item if $\gA=\gD=\gE=0$ and $\lambda_2>0$, the unique solution is the
    ridge-GMV portfolio
    $\w^\star=(\Sighat+\lambda_2I)^{-1}\one/\big(\one^\top(\Sighat+\lambda_2I)^{-1}\one\big)$.
\end{enumerate}
\end{lemma}

\begin{proof}
(ii) implies (i) on setting $\lambda_2=0$, so we prove (ii). With only the budget
constraint, \eqref{eq:kkt} reduces to stationarity of
$L(\w,\nu)=\w^\top(\Sighat+\lambda_2I)\w-\nu(\one^\top\w-1)$, namely
$2(\Sighat+\lambda_2I)\w^\star=\nu\one$. Since $\Sighat+\lambda_2I\succ0$ it is
invertible, so $\w^\star=\tfrac{\nu}{2}(\Sighat+\lambda_2I)^{-1}\one$. Imposing
$\one^\top\w^\star=1$ gives
$\tfrac{\nu}{2}=\big(\one^\top(\Sighat+\lambda_2I)^{-1}\one\big)^{-1}$, which is
well defined and positive because $(\Sighat+\lambda_2I)^{-1}\succ0$. Substituting
yields the stated formula. Uniqueness follows from
Theorem~\ref{thm:uniqueness}(i)--(ii).
\end{proof}

Lemma~\ref{lem:reductions} is what makes the formulation a strict generalization
rather than a substitute: the classical portfolios are recovered exactly, not
approximately, and the two identities serve as numerical acceptance tests for the
implementation (Section~\ref{sec:computation}).

\subsection{Monotone comparative statics}

The next result explains, and indeed predicts, the monotone behaviour observed in
the penalty sweeps of Section~\ref{sec:computation}. Write the objective as
$F_\gamma=f_0+\gamma h$, where $h$ is one of the design criteria and $f_0$
collects all remaining terms.

\begin{theorem}[Monotonicity in the penalty strength]
\label{thm:monotone}
Let $f_0$ and $h$ be convex on the convex set $\Wset\cap\dom F$, and for
$\gamma\ge0$ let $\w_\gamma$ be any minimizer of $f_0+\gamma h$ over that set.
Then for $0\le\gamma<\gamma'$,
\[
  h(\w_{\gamma'})\;\le\;h(\w_{\gamma})
  \qquad\text{and}\qquad
  f_0(\w_{\gamma'})\;\ge\;f_0(\w_{\gamma}).
\]
That is, increasing a penalty weakly improves its own criterion and weakly
worsens every other component of the objective.
\end{theorem}

\begin{proof}
Optimality of $\w_\gamma$ at $\gamma$ and of $\w_{\gamma'}$ at $\gamma'$ give
\[
  f_0(\w_\gamma)+\gamma h(\w_\gamma)\le f_0(\w_{\gamma'})+\gamma h(\w_{\gamma'}),
  \qquad
  f_0(\w_{\gamma'})+\gamma' h(\w_{\gamma'})\le f_0(\w_\gamma)+\gamma' h(\w_\gamma).
\]
Adding the two inequalities and cancelling $f_0(\w_\gamma)+f_0(\w_{\gamma'})$
yields $\gamma h(\w_\gamma)+\gamma'h(\w_{\gamma'})\le\gamma h(\w_{\gamma'})+\gamma'h(\w_\gamma)$,
i.e.\ $(\gamma'-\gamma)\big(h(\w_\gamma)-h(\w_{\gamma'})\big)\ge0$. Since
$\gamma'>\gamma$ this gives $h(\w_{\gamma'})\le h(\w_\gamma)$. Substituting this
into the first displayed inequality,
$f_0(\w_\gamma)-f_0(\w_{\gamma'})\le\gamma\big(h(\w_{\gamma'})-h(\w_\gamma)\big)\le0$,
which is the second claim.
\end{proof}

\begin{remark}
Theorem~\ref{thm:monotone} is the formal counterpart of two empirical
observations reported below: that each criterion sweep moves its own target
metric monotonically (Figure~\ref{fig:sweep}), and that the A-penalized portfolio
improves its design criterion while giving up portfolio-variance performance. The
second is not a defect of the A-criterion but an instance of the general
trade-off the theorem describes.
\end{remark}

\subsection{Limiting solutions}

\begin{lemma}[Penalty limit]
\label{lem:penaltylimit}
Let $\Wset$ be compact, let $f_0$ be bounded below on $\Wset$ with
$f_0(\bar\w)<\infty$, and let $h$ be closed convex with a \emph{unique} minimizer
$\bar\w$ over $\Wset$. If $\w_\gamma$ minimizes $f_0+\gamma h$ over $\Wset$, then
$\w_\gamma\to\bar\w$ as $\gamma\to\infty$.
\end{lemma}

\begin{proof}
Optimality of $\w_\gamma$ gives
$f_0(\w_\gamma)+\gamma h(\w_\gamma)\le f_0(\bar\w)+\gamma h(\bar\w)$, hence
\[
  0\le h(\w_\gamma)-h(\bar\w)\le\frac{f_0(\bar\w)-f_0(\w_\gamma)}{\gamma}
  \le\frac{f_0(\bar\w)-\inf_{\Wset}f_0}{\gamma}\xrightarrow[\gamma\to\infty]{}0 ,
\]
the first inequality because $\bar\w$ minimizes $h$. Thus
$h(\w_\gamma)\to h(\bar\w)=\min_{\Wset}h$. Let $\w'$ be any accumulation point of
$(\w_\gamma)$, which exists by compactness. Lower semicontinuity of $h$ gives
$h(\w')\le\liminf h(\w_\gamma)=\min_{\Wset}h$, so $\w'$ minimizes $h$ and
therefore $\w'=\bar\w$ by uniqueness. As every accumulation point equals
$\bar\w$ and $\Wset$ is compact, $\w_\gamma\to\bar\w$.
\end{proof}

\begin{theorem}[Limiting portfolios]
\label{thm:limits}
Let $\Wset=\{\w:\one^\top\w=1,\ \w\ge0\}$ be the unit simplex and $K=N$. Then:
\begin{enumerate}[label=(\roman*),leftmargin=2.0em,itemsep=0.2em]
  \item \emph{(D-limit)} as $\gD\to\infty$ with the other penalties fixed,
    $\w^\star\to\tfrac1N\one$, the equally weighted portfolio;
  \item \emph{(A-limit)} as $\gA\to\infty$ with the other penalties fixed,
    $\w^\star\to\big(\sqrt{G_{ii}}\big/\sum_{j}\sqrt{G_{jj}}\big)_{i=1}^{N}$,
    where $G=(V^\top V)^{-1}$;
  \item \emph{(ridge limit)} as $\lambda_2\to\infty$ with the other penalties
    fixed, $\w^\star\to\tfrac1N\one$;
  \item \emph{(E-limit)} as $\gE\to\infty$, every accumulation point of
    $\w^\star$ maximizes $\lammin(\M)$ over $\Wset$, i.e.\ is an E-optimal
    design; if that maximizer is unique, $\w^\star$ converges to it.
\end{enumerate}
\end{theorem}

\begin{proof}
All four statements apply Lemma~\ref{lem:penaltylimit}; it remains to identify the
minimizer of each penalty over the simplex.

(i) By Corollary~\ref{cor:KN} the D-criterion is $-\sum_i\log w_i$ up to a
constant. It is strictly convex on the open simplex, so its minimizer is unique
and characterized by stationarity of
$-\sum_i\log w_i-\nu(\one^\top\w-1)$: $-1/w_i=\nu$ for all $i$, so all $w_i$ are
equal, giving $\w=\tfrac1N\one$. (Equivalently, by the AM--GM inequality
$\prod_iw_i\le(1/N)^N$ on the simplex with equality iff $\w=\tfrac1N\one$.)

(ii) By Corollary~\ref{cor:KN} the A-criterion is $\sum_iG_{ii}/w_i$, strictly
convex on the open simplex. Stationarity gives $-G_{ii}/w_i^2=\nu$, so
$w_i\propto\sqrt{G_{ii}}$, and normalizing by the budget yields the stated
vector. (Equivalently, by Cauchy--Schwarz,
$\big(\sum_i\sqrt{G_{ii}}\big)^2\le\big(\sum_iG_{ii}/w_i\big)\big(\sum_iw_i\big)$
with equality iff $w_i\propto\sqrt{G_{ii}}$.)

(iii) $\lVert\w\rVert_2^2$ is strictly convex with unique simplex minimizer
$\tfrac1N\one$, by the same stationarity argument or by
$\lVert\w\rVert_2^2\ge(\one^\top\w)^2/N=1/N$ with equality iff $\w=\tfrac1N\one$.

(iv) The E-criterion $-\lammin(\M)$ is convex and continuous on the compact
simplex, so minimizers exist, but it need not be strictly convex and the
minimizer need not be unique. The proof of Lemma~\ref{lem:penaltylimit} up to the
accumulation-point step applies verbatim and shows that every accumulation point
minimizes $-\lammin(\M)$, i.e.\ maximizes $\lammin(\M)$; uniqueness then upgrades
this to convergence.
\end{proof}

\begin{remark}[Two routes to $1/N$]
Parts (i) and (iii) both terminate at the equally weighted portfolio, but by
different mechanisms and along different paths: the ridge penalizes \emph{large}
weights quadratically, whereas the D-barrier penalizes \emph{small} weights
without bound. The D-path therefore approaches $1/N$ from the interior while
enforcing strict positivity at every finite $\gD$ (Proposition~\ref{prop:interior}),
which the ridge does not. This distinction is invisible in the limit but governs
the behaviour at the moderate penalty values used in practice.
\end{remark}

\subsection{Stability: why the penalties reduce turnover}

The empirical headline of this paper is a large reduction in turnover. Turnover
is a difference of consecutive solutions, so it is controlled by the modulus of
continuity of the solution map $\Sighat\mapsto\w^\star(\Sighat)$. We now show that
this map is Lipschitz with a constant inversely proportional to the
strong-convexity modulus, and that the penalties increase that modulus.

\begin{lemma}[Strong-convexity modulus]
\label{lem:modulus}
Let $K=N$, let $\Wset$ be contained in the unit simplex, and let
$G=(V^\top V)^{-1}$. Then $F$ is $\mu$-strongly convex on $\Wset\cap\dom F$ with
\begin{equation}
  \mu\;\ge\;2\lammin\!\big(\Sighat\big)+2\lambda_2+\gD+2\gA\min_iG_{ii} .
  \label{eq:mu}
\end{equation}
\end{lemma}

\begin{proof}
The Hessians of the risk and ridge terms are $2\Sighat\succeq2\lammin(\Sighat)I$
and $2\lambda_2I$. By Corollary~\ref{cor:KN} the D- and A-terms are
$-\gD\sum_i\log w_i$ and $\gA\sum_iG_{ii}/w_i$ with Hessians
$\gD\diag(w_i^{-2})$ and $2\gA\diag(G_{ii}w_i^{-3})$. On the unit simplex
$0<w_i\le1$, so $w_i^{-2}\ge1$ and $w_i^{-3}\ge1$, giving
$\gD\diag(w_i^{-2})\succeq\gD I$ and
$2\gA\diag(G_{ii}w_i^{-3})\succeq2\gA(\min_iG_{ii})I$. The E-term is convex, so
contributes a positive semidefinite term in the subdifferential sense. Summing
the lower bounds gives \eqref{eq:mu}.
\end{proof}

\begin{theorem}[Stability of the solution map]
\label{thm:stability}
Let $\Wset$ be convex and compact, and let $F(\cdot;\Sighat)$ be $\mu$-strongly
convex on $\Wset\cap\dom F$ uniformly in the risk model, with minimizers
$\w^\star(\Sighat)$. Then for any two covariance estimates $\Sighat_1,\Sighat_2$,
\begin{equation}
  \big\lVert\w^\star(\Sighat_1)-\w^\star(\Sighat_2)\big\rVert_2
  \;\le\;\frac{2}{\mu}\,\big\lVert\Sighat_1-\Sighat_2\big\rVert_2\,
  \big\lVert\w^\star(\Sighat_2)\big\rVert_2 .
  \label{eq:stability}
\end{equation}
In particular, if $\Wset$ lies in the unit simplex then
$\lVert\w^\star(\Sighat_2)\rVert_2\le1$ and the solution map is Lipschitz with
constant $2/\mu$.
\end{theorem}

\begin{proof}
Write $F_k(\w)=F(\w;\Sighat_k)$ and $\w_k=\w^\star(\Sighat_k)$ for $k=1,2$, and
put $\Delta=\Sighat_2-\Sighat_1$. Since the two objectives differ only in the
risk term, $F_2(\w)-F_1(\w)=\w^\top\Delta\w$ and hence
\begin{equation}
  \nabla F_2(\w)=\nabla F_1(\w)+2\Delta\w .
  \label{eq:gradshift}
\end{equation}
First-order optimality for a convex function on a convex set states
$\langle\nabla F_k(\w_k),\w-\w_k\rangle\ge0$ for all $\w\in\Wset$. Taking
$\w=\w_2$ in the case $k=1$ and $\w=\w_1$ in the case $k=2$ and adding,
\[
  \big\langle\nabla F_1(\w_1)-\nabla F_2(\w_2),\,\w_2-\w_1\big\rangle\ \ge\ 0 .
\]
Substituting \eqref{eq:gradshift} at $\w=\w_2$ and rearranging,
\[
  \big\langle\nabla F_1(\w_1)-\nabla F_1(\w_2),\,\w_2-\w_1\big\rangle
  \;\ge\;2\big\langle\Delta\w_2,\,\w_2-\w_1\big\rangle .
\]
By $\mu$-strong convexity of $F_1$, the gradient is strongly monotone,
$\langle\nabla F_1(\w_2)-\nabla F_1(\w_1),\w_2-\w_1\rangle\ge\mu\lVert\w_2-\w_1\rVert_2^2$,
so the left-hand side is at most $-\mu\lVert\w_2-\w_1\rVert_2^2$. Therefore
\[
  \mu\lVert\w_2-\w_1\rVert_2^2
  \;\le\;-2\big\langle\Delta\w_2,\,\w_2-\w_1\big\rangle
  \;\le\;2\lVert\Delta\w_2\rVert_2\,\lVert\w_2-\w_1\rVert_2
  \;\le\;2\lVert\Delta\rVert_2\lVert\w_2\rVert_2\lVert\w_2-\w_1\rVert_2 ,
\]
using Cauchy--Schwarz and the operator-norm bound. Dividing by
$\mu\lVert\w_2-\w_1\rVert_2$ (the claim being trivial if $\w_1=\w_2$) gives
\eqref{eq:stability}.
\end{proof}

\begin{corollary}[Turnover bound]
\label{cor:turnover}
Consider a rebalancing scheme that at date $t$ solves \eqref{eq:objective} with
covariance estimate $\Sighat_t$ over a fixed feasible set $\Wset$ contained in the
unit simplex. Then the realized turnover satisfies
\begin{equation}
  \tau_t\;=\;\big\lVert\w^\star_t-\w^\star_{t-1}\big\rVert_1
  \;\le\;\sqrt{N}\,\big\lVert\w^\star_t-\w^\star_{t-1}\big\rVert_2
  \;\le\;\frac{2\sqrt{N}}{\mu}\,\big\lVert\Sighat_t-\Sighat_{t-1}\big\rVert_2 ,
\end{equation}
with $\mu$ bounded below as in \eqref{eq:mu}. The bound is decreasing in each of
$\lambda_2$, $\gD$ and $\gA$.
\end{corollary}

\begin{proof}
The first inequality is the norm equivalence
$\lVert x\rVert_1\le\sqrt{N}\lVert x\rVert_2$ on $\R^{N}$; the second is
Theorem~\ref{thm:stability} with $\lVert\w^\star_{t-1}\rVert_2\le1$. Monotonicity
in the penalties is immediate from \eqref{eq:mu}, since $\mu$ appears in the
denominator and its lower bound is increasing in $\lambda_2$, $\gD$ and $\gA$.
\end{proof}

Corollary~\ref{cor:turnover} is, to our knowledge, the first result that derives
the turnover-damping effect of this class of regularizers from the geometry of
the objective rather than observing it empirically. It says that turnover is
bounded by the instability of the covariance estimate divided by the
strong-convexity modulus, and that every penalty in \eqref{eq:objective} except
the E-term enters that modulus additively. The empirical turnover reductions
reported in Sections~\ref{sec:computation}--\ref{sec:applications}---roughly an
order of magnitude relative to sample-covariance Markowitz---are consistent with
this mechanism.

\begin{remark}[Where the E-term sits]
The bound \eqref{eq:mu} credits no strong convexity to the E-term, because
$-\lammin(\M)$ is convex but not strongly convex: it is a pointwise minimum of
linear functions and is affine along directions that do not perturb the minimal
eigenvalue. The E-penalty therefore stabilizes the portfolio through a different
channel---raising $\lammin(\M)$ and lowering the condition number, which improves
the conditioning of the problem rather than its curvature. This is consistent
with our empirical finding that the E-variant tracks GMV more closely than the
D-variant does on turnover.
\end{remark}
